% ================= 附录 A: 命题 2.1.4 的证明 =================

\appendix
\section{命题 2.1.4 的证明}

\subsection{一些拓扑预备}\label{sec:A.1}

设 $\bigl(\pind{}_{B(\Qp)}^{\GL_2(\Qp)} \chi_1 \otimes \chi_2\bigr)(\Qp)^{\mathrm{an}}$ (resp.\ $\bigl(\pind{}_{B(\Qp)}^{\GL_2(\Qp)} \chi_1 \otimes \chi_2\bigr)(\Zp)^{\mathrm{an}}$) 为 $\anind_{B(\Qp)}^{\GL_2(\Qp)} \chi_1 \otimes \chi_2$ 中支撑在 $\Qp$ 内紧 (resp.\ 支撑在 $\Zp$ 内) 的局部解析函数 $f : \Qp \to L$ (参见 \ref{sec:2.1}) 组成的闭子空间.
它被 $(\mathfrak{gl}_2(\Qp),B(\Qp))$ (resp.\ $(\mathfrak{gl}_2(\Qp),B(\Qp)^{+})$) 稳定, 且自然映射:
\begin{equation}
\label{eq:35}
L[N(\Qp)] \otimes_{L[N(\Zp)]} \bigl(\pind{}_{B(\Qp)}^{\GL_2(\Qp)} \chi_1 \otimes \chi_2\bigr)(\Zp)^{\mathrm{an}}
\to \bigl(\pind{}_{B(\Qp)}^{\GL_2(\Qp)} \chi_1 \otimes \chi_2\bigr)(\Qp)^{\mathrm{an}}
\end{equation}
是同构.
还记 $\bigl(\pind{}_{B(\Qp)}^{\GL_2(\Qp)} \chi_1 \otimes \chi_2\bigr)(\Zp)^{\mathrm{lp}}$ 和 $\bigl(\pind{}_{B(\Qp)}^{\GL_2(\Qp)} \chi_1 \otimes \chi_2\bigr)(\Zp)^{\mathrm{lp},\le d}$ 为 $\bigl(\pind{}_{B(\Qp)}^{\GL_2(\Qp)} \chi_1 \otimes \chi_2\bigr)(\Qp)^{\mathrm{lp}}$ 和 $\bigl(\pind{}_{B(\Qp)}^{\GL_2(\Qp)} \chi_1 \otimes \chi_2\bigr)(\Qp)^{\mathrm{lp},\le d}$ 中支撑在 $\Zp$ 内的函数组成的闭子空间.
它们仍是被 $(\mathfrak{gl}_2(\Qp),B(\Qp)^{+})$ 稳定的子空间, 且与 \eqref{eq:35} 类似的映射是同构.

若 $r \ge 1$, 记 $\bigl(\pind{}_{B(\Qp)}^{\GL_2(\Qp)} \chi_1 \otimes \chi_2\bigr)(\Zp)^{\mathrm{an}}_{r}$ 为 $\Zp$ 上局部解析、且在 $\Zp$ 的每个闭圆盘 $D(z,r) := \{x \in \Qp \mid \|x-z\| \le |p^{r}|\}$ ($z \in \Zp$) 上解析的函数组成的 Banach 空间.
若 $\mathcal{O}(D(z,r))$ 表示 $D(z,r)$ 上取值于 $L$ 的刚性解析函数的 Banach 空间, 有自然拓扑同构:
\[
\bigl(\pind{}_{B(\Qp)}^{\GL_2(\Qp)} \chi_1 \otimes \chi_2\bigr)(\Zp)^{\mathrm{an}}_{r}
\xrightarrow{\sim} \bigoplus_{z \in \Zp/p^{r}} \mathcal{O}(D(z,r))
\]
(注意: $\bigl(\pind{}_{B(\Qp)}^{\GL_2(\Qp)} \chi_1 \otimes \chi_2\bigr)(\Zp)^{\mathrm{an}}_{r}$ 由 $\bigl(\pind{}_{B(\Qp)}^{\GL_2(\Qp)} \chi_1 \otimes \chi_2\bigr)(\Zp)^{\mathrm{an}}$ 诱导的拓扑比上述 Banach 空间拓扑粗糙).
若选择 $r_0 \ge 1$ 使得 $\chi_2\chi_1^{-1}$ 在 $1 + p^{r_0}\Zp$ 上解析, 则对一切 $r \ge r_0$, $\bigl(\pind{}_{B(\Qp)}^{\GL_2(\Qp)} \chi_1 \otimes \chi_2\bigr)(\Zp)^{\mathrm{an}}_{r}$ 在 $\bigl(\pind{}_{B(\Qp)}^{\GL_2(\Qp)} \chi_1 \otimes \chi_2\bigr)(\Zp)^{\mathrm{an}}$ 中被 $K_p(r)$ 稳定 (关于 $K_p(r)$ 参见 \ref{sec:3.1}).
此外容易验证 $K_p(r)$ 的作用对 $\bigl(\pind{}_{B(\Qp)}^{\GL_2(\Qp)} \chi_1 \otimes \chi_2\bigr)(\Zp)^{\mathrm{an}}_{r}$ 的 Banach 空间拓扑是连续的.
事实上还有更多: 若 $\mathcal{K}_p(r)$ 表示 $\Qp$ 上代数群 $\GL_2$ 的由条件 $|a-1|,|b|,|c|,|d-1| \le |p^{r}|$ 定义的仿射型子群, 则 $K_p(r)$ 是 $\mathcal{K}_p(r)$ 的 $\Qp$-点群, 且 $K_p(r)$ 的作用实际上由 $\mathcal{K}_p(r)$ 在 $\bigl(\pind{}_{B(\Qp)}^{\GL_2(\Qp)} \chi_1 \otimes \chi_2\bigr)(\Zp)^{\mathrm{an}}_{r}$ 上的刚性解析作用诱导 ($K_p(r)$ 的作用按 \cite{ref:26} 的意义是"$\mathcal{K}_p(r)$-解析的").
特别地, 可以微分这个作用, 得到 $\mathfrak{gl}_2(\Qp)$ 在 $\bigl(\pind{}_{B(\Qp)}^{\GL_2(\Qp)} \chi_1 \otimes \chi_2\bigr)(\Zp)^{\mathrm{an}}_{r}$ 上的作用 (对 $r \ge r_0$).
自然映射:
\begin{equation}
\label{eq:36}
\varinjlim_{r \ge r_0} \bigl(\pind{}_{B(\Qp)}^{\GL_2(\Qp)} \chi_1 \otimes \chi_2\bigr)(\Zp)^{\mathrm{an}}_{r}
\to \bigl(\pind{}_{B(\Qp)}^{\GL_2(\Qp)} \chi_1 \otimes \chi_2\bigr)(\Zp)^{\mathrm{an}}
\end{equation}
于是是一个 $\mathfrak{gl}_2(\Qp)$-等变拓扑同构.

对 $d \ge 0$ 和 $r \ge 1$, 设:
\[
\bigl(\pind{}_{B(\Qp)}^{\GL_2(\Qp)} \chi_1 \otimes \chi_2\bigr)(\Zp)^{\mathrm{lp},\le d}_{r}
:= \bigl(\pind{}_{B(\Qp)}^{\GL_2(\Qp)} \chi_1 \otimes \chi_2\bigr)(\Zp)^{\mathrm{an}}_{r}
\cap \bigl(\pind{}_{B(\Qp)}^{\GL_2(\Qp)} \chi_1 \otimes \chi_2\bigr)(\Zp)^{\mathrm{lp},\le d},
\]
以及:
\[
\bigl(\pind{}_{B(\Qp)}^{\GL_2(\Qp)} \chi_1 \otimes \chi_2\bigr)(\Zp)^{\mathrm{lp}}_{r}
:= \bigl(\pind{}_{B(\Qp)}^{\GL_2(\Qp)} \chi_1 \otimes \chi_2\bigr)(\Zp)^{\mathrm{an}}_{r}
\cap \bigl(\pind{}_{B(\Qp)}^{\GL_2(\Qp)} \chi_1 \otimes \chi_2\bigr)(\Zp)^{\mathrm{lp}}
= \bigcup_{d \ge 0} \bigl(\pind{}_{B(\Qp)}^{\GL_2(\Qp)} \chi_1 \otimes \chi_2\bigr)(\Zp)^{\mathrm{lp},\le d}_{r}.
\]
每个空间 $\bigl(\pind{}_{B(\Qp)}^{\GL_2(\Qp)} \chi_1 \otimes \chi_2\bigr)(\Zp)^{\mathrm{lp},\le d}_{r}$ 都是有限维的, 其并对 $\bigl(\pind{}_{B(\Qp)}^{\GL_2(\Qp)} \chi_1 \otimes \chi_2\bigr)(\Zp)^{\mathrm{an}}_{r}$ 的 (Banach) 拓扑稠密.
若 $d$ 足够大, 则 $\bigl(\pind{}_{B(\Qp)}^{\GL_2(\Qp)} \chi_1 \otimes \chi_2\bigr)(\Zp)^{\mathrm{lp},\le d}_{r}$ 作为 $\mathfrak{gl}_2(\Qp)$-模生成 $\bigl(\pind{}_{B(\Qp)}^{\GL_2(\Qp)} \chi_1 \otimes \chi_2\bigr)(\Zp)^{\mathrm{lp}}_{r}$ (事实上, 若 $\chi_1 \otimes \chi_2$ 不是经典权的, 任何 $d \ge 0$ 都可以; 但若 $\chi_1 \otimes \chi_2$ 是经典权 $k \ge 2$ 的, 需取 $d \ge k-1$, 参见引理 \ref{lem:2.1.2} 和引理 \ref{lem:2.1.3}).

现在回到命题 \ref{prop:2.1.4} 的 $\GL_2(\Qp)$-Banach 空间 $V$.
对 $r \ge 1$, 设 $V^{K_p(r)\text{-an}}$ 为 $V$ 中满足如下条件的向量 $v$ 的子空间: "轨道"映射 $g \to g \cdot v$ 限制到 $K_p(r)$ 时是 $K_p(r)$ 上取值于 $V$ 的 (某个 $K_p(r)$ 上的) 刚性解析函数.
若 $\mathcal{O}(K_p(r),V)$ 是 $K_p(r)$ 上取值于 $V$ 的刚性解析函数的 Banach 空间, 则把 $v \in V^{K_p(r)\text{-an}}$ 送到其轨道映射的箭头把 $V^{K_p(r)\text{-an}}$ 等同于 $\mathcal{O}(K_p(r),V)$ 的一个闭子空间 (\cite[Prop.~3.3.3]{ref:26}).
给 $V^{K_p(r)\text{-an}}$ 配备由 $\mathcal{O}(K_p(r),V)$ 的拓扑诱导的 Banach 拓扑 (它比由 Banach 空间 $V$ 诱导的拓扑细).
则有拓扑同构:
\[
\varinjlim_r V^{K_p(r)\text{-an}} \xrightarrow{\sim} V^{\mathrm{an}}.
\]

\begin{remark}\label{rem:A.1.1}
这里使用 \cite[Def.~3.5.3]{ref:26} 中给出的局部解析向量定义, 它与 \cite{ref:45} 的定义不同: 两个定义给出相同的底层空间 $V^{\mathrm{an}}$, 但 \cite{ref:26} 定义的 $V^{\mathrm{an}}$ 的拓扑一般比 \cite{ref:45} 定义的细 (参见 \cite[Th.~3.5.7]{ref:26}).
我们将在 $V^{\mathrm{an}}$ 配备 \cite{ref:26} 拓扑的情形证明命题 \ref{prop:2.1.4}, 这立即蕴含 $V^{\mathrm{an}}$ 配备 \cite{ref:45} 的更粗拓扑时的相应陈述.
注意在应用中, $V$ 总是可容许的 (酉) $\GL_2(\Qp)$-Banach 空间, 此时 $V^{\mathrm{an}}$ 上的两种拓扑相同 (利用 \cite[Prop.~6.2.4]{ref:26} 和 \cite[Th.~7.1]{ref:45}).
\end{remark}

\subsection{命题 2.1.4 的证明}

保留 \ref{sec:A.1} 的记号.

\begin{lemma}\label{lem:A.2.1}
限制:
\begin{equation}
\label{eq:37}
\Hom_{B(\Qp)^{+}}\Bigl(\bigl(\pind{}_{B(\Qp)}^{\GL_2(\Qp)} \chi_1 \otimes \chi_2\bigr)(\Zp)^{\mathrm{an}},\ V\Bigr)
\to \Hom_{B(\Qp)^{+}}\Bigl(\bigl(\pind{}_{B(\Qp)}^{\GL_2(\Qp)} \chi_1 \otimes \chi_2\bigr)(\Zp)^{\mathrm{lp}},\ V\Bigr)
\end{equation}
是同构.
\end{lemma}

\begin{proof}
由于 $\bigl(\pind{}_{B(\Qp)}^{\GL_2(\Qp)} \chi_1 \otimes \chi_2\bigr)(\Zp)^{\mathrm{lp}}$ 在空间 $\bigl(\pind{}_{B(\Qp)}^{\GL_2(\Qp)} \chi_1 \otimes \chi_2\bigr)(\Zp)^{\mathrm{an}}$ 中稠密, 连续线性映射:
\begin{equation}
\label{eq:38}
\bigl(\pind{}_{B(\Qp)}^{\GL_2(\Qp)} \chi_1 \otimes \chi_2\bigr)(\Zp)^{\mathrm{an}} \to V
\end{equation}
为零 (resp.\ 是 $B(\Qp)^{+}$-等变的) 当且仅当它限制到 $\bigl(\pind{}_{B(\Qp)}^{\GL_2(\Qp)} \chi_1 \otimes \chi_2\bigr)(\Zp)^{\mathrm{lp}}$ 上如此.
因此只需证明 \eqref{eq:37} 像中的任何元素都能延拓为连续线性映射 \eqref{eq:38}.
若 $f(X) \in \Qp[X]$, 记 $\|f(X)\|$ 为其系数的 $p$ 进绝对值的上确界.
记 $\|\cdot\|_V$ 为 $V$ 上给出其 Banach 拓扑的一个范数.
选择 $\|\cdot\|_V$ 在 $N(\Zp)$ 下不变 (这是可能的, 因为 $N(\Zp)$ 紧).
设 $A > 0$ 使得对一切 $v \in V$ 有 $\|\wp v\|_V \le A\|v\|_V$, 其中 $\wp := \begin{pmatrix} p & 0 \\ 0 & 1 \end{pmatrix}$.
设 $\varphi$ 在 \eqref{eq:37} 的像中, 对每个 $d \ge 0$, 设 $C(d) > 0$ 使得对次数至多 $d$ 的 $f(X)$ 有 $\|\varphi(f(X))\|_V \le C(d)\|f(X)\|$.
对每个 $z \in \Zp$ 和 $r \ge 0$, 有 (参见 \eqref{eq:1}):
\[
f\!\left(\frac{X-z}{p^{r}}\right)\Big|_{z+p^{r}\Zp}
= \chi_2(p)^{r}\begin{pmatrix} 1 & z \\ 0 & 1 \end{pmatrix}\wp^{r} \cdot f(X).
\]
利用 $\varphi$ 在 $B(\Qp)^{+}$ 下的不变性, 得到:
\[
\Bigl\|\varphi\Bigl(f\!\Bigl(\frac{X-z}{p^{r}}\Bigr)\Bigr)\Bigr\|_V
= |\chi_2(p)|^{-r}\Bigl\|\begin{pmatrix} 1 & z \\ 0 & 1 \end{pmatrix}\wp^{r} \cdot \varphi(f(X))\Bigr\|_V
\le \bigl(|\chi_2(p)|^{-1}A\bigr)^{r}\|\varphi(f(X))\|_V
\le \bigl(|\chi_2(p)|^{-1}A\bigr)^{r}C(d)\|f(X)\|.
\]
取 $h$ 使得 $|\chi_2(p)|^{-1}A \le |p^{-h}|$, 可见对每个 $d \ge 0$, $\varphi$ 限制到 $\bigl(\pind{}_{B(\Qp)}^{\GL_2(\Qp)} \chi_1 \otimes \chi_2\bigr)(\Zp)^{\mathrm{lp},\le d}$ 上是阶 $\le h$ 的缓增映射.
由 Amice--V\'elu (\cite{ref:1}) 和 Vishik (\cite{ref:47}) 的定理, 若 $d > h-1$, 则限制唯一地延拓为连续线性映射 $\widetilde{\varphi}_d : \bigl(\pind{}_{B(\Qp)}^{\GL_2(\Qp)} \chi_1 \otimes \chi_2\bigr)(\Zp)^{\mathrm{an}} \to V$, 它是阶 $\le h$ 的缓增映射.
唯一性表明延拓 $\widetilde{\varphi}_d$ 不依赖于 $d > h-1$ 的选择, 并给出 $\varphi$ 到 $\bigl(\pind{}_{B(\Qp)}^{\GL_2(\Qp)} \chi_1 \otimes \chi_2\bigr)(\Zp)^{\mathrm{an}}$ 的所求延拓.
\end{proof}

\begin{lemma}\label{lem:A.2.2}
限制:
\[
\Hom_{(\mathfrak{gl}_2(\Qp),B(\Qp))}\Bigl(\bigl(\pind{}_{B(\Qp)}^{\GL_2(\Qp)} \chi_1 \otimes \chi_2\bigr)(\Qp)^{\mathrm{an}},\ V^{\mathrm{an}}\Bigr)
\to \Hom_{(\mathfrak{gl}_2(\Qp),B(\Qp))}\Bigl(\bigl(\pind{}_{B(\Qp)}^{\GL_2(\Qp)} \chi_1 \otimes \chi_2\bigr)(\Qp)^{\mathrm{lp}},\ V^{\mathrm{an}}\Bigr)
\]
是同构.
\end{lemma}

\begin{proof}
由于 \eqref{eq:35} 及其把"an"换成"lp"的类似陈述都是同构, 且 $B(\Qp)^{+}$ 生成群 $B(\Qp)$, 只需证明限制:
\begin{equation}
\label{eq:39}
\Hom_{(\mathfrak{gl}_2(\Qp),B(\Qp)^{+})}\Bigl(\bigl(\pind{}_{B(\Qp)}^{\GL_2(\Qp)} \chi_1 \otimes \chi_2\bigr)(\Zp)^{\mathrm{an}},\ V^{\mathrm{an}}\Bigr)
\to \Hom_{(\mathfrak{gl}_2(\Qp),B(\Qp)^{+})}\Bigl(\bigl(\pind{}_{B(\Qp)}^{\GL_2(\Qp)} \chi_1 \otimes \chi_2\bigr)(\Zp)^{\mathrm{lp}},\ V^{\mathrm{an}}\Bigr)
\end{equation}
是同构.
考虑到引理 \ref{lem:A.2.1}, 只需证明: 若 $\varphi : \bigl(\pind{}_{B(\Qp)}^{\GL_2(\Qp)} \chi_1 \otimes \chi_2\bigr)(\Zp)^{\mathrm{an}} \to V$ 是连续映射, 其限制 $\varphi_{\mathrm{lp}}$ 到 $\bigl(\pind{}_{B(\Qp)}^{\GL_2(\Qp)} \chi_1 \otimes \chi_2\bigr)(\Zp)^{\mathrm{lp}}$ 能分解为连续 $\mathfrak{gl}_2(\Qp)$-等变映射 $\bigl(\pind{}_{B(\Qp)}^{\GL_2(\Qp)} \chi_1 \otimes \chi_2\bigr)(\Zp)^{\mathrm{lp}} \to V^{\mathrm{an}}$ 与自然注入 $V^{\mathrm{an}} \to V$ 的复合, 则 $\varphi$ 也有类似的分解.
取 $r \ge r_0$, 设 $\varphi_r$ 为 $\varphi$ 限制到 $\bigl(\pind{}_{B(\Qp)}^{\GL_2(\Qp)} \chi_1 \otimes \chi_2\bigr)(\Zp)^{\mathrm{an}}_{r}$, 并取 $d$ 足够大使得 $\bigl(\pind{}_{B(\Qp)}^{\GL_2(\Qp)} \chi_1 \otimes \chi_2\bigr)(\Zp)^{\mathrm{lp},\le d}_{r}$ 在 $\mathfrak{gl}_2(\Qp)$ 作用下生成 $\bigl(\pind{}_{B(\Qp)}^{\GL_2(\Qp)} \chi_1 \otimes \chi_2\bigr)(\Zp)^{\mathrm{an}}_{r}$.
由于 $\bigl(\pind{}_{B(\Qp)}^{\GL_2(\Qp)} \chi_1 \otimes \chi_2\bigr)(\Zp)^{\mathrm{lp},\le d}_{r}$ 有限维, 其经 $\varphi_{\mathrm{lp}}$ 在 $V^{\mathrm{an}}$ 中的像有限维, 因而落在某个 $s \ge r$ 的 $V^{K_p(s)\text{-an}}$ 中.
由于 $\varphi_{\mathrm{lp}}$ 是 $\mathfrak{gl}_2(\Qp)$-等变的且 $V^{K_p(s)\text{-an}}$ 被 $\mathfrak{gl}_2(\Qp)$ 稳定, 可见 $\bigl(\pind{}_{B(\Qp)}^{\GL_2(\Qp)} \chi_1 \otimes \chi_2\bigr)(\Zp)^{\mathrm{lp}}_{r}$ 的像仍落在 $V^{K_p(s)\text{-an}}$ 中.
现在考虑连续映射:
\[
\Phi : K_p(s) \times \bigl(\pind{}_{B(\Qp)}^{\GL_2(\Qp)} \chi_1 \otimes \chi_2\bigr)(\Zp)^{\mathrm{an}}_{r} \to V,
\qquad (g,f) \mapsto \varphi_r(g \cdot f) - g \cdot \varphi_r(f).
\]
设 $f \in \bigl(\pind{}_{B(\Qp)}^{\GL_2(\Qp)} \chi_1 \otimes \chi_2\bigr)(\Zp)^{\mathrm{lp}}_{r}$, 则映射 $g \mapsto \Phi(g,f)$ 是 $K_p(s)$ 上 ($K_p(s)$ 上某刚性解析函数的限制得到的) 刚性解析函数 (因为 $K_p(s)$ 对 $f$ 和 $\varphi_r(f)$ 的作用来自 $K_p(s)$ 的刚性解析作用).
此外, 由于 $\varphi_{\mathrm{lp}}$ 是 $\mathfrak{gl}_2(\Qp)$-等变的, 它的所有导数都为零.
因此 $\Phi$ 在 $K_p(s) \times \bigl(\pind{}_{B(\Qp)}^{\GL_2(\Qp)} \chi_1 \otimes \chi_2\bigr)(\Zp)^{\mathrm{lp}}_{r}$ 上为零.
但这个空间在 $K_p(s) \times \bigl(\pind{}_{B(\Qp)}^{\GL_2(\Qp)} \chi_1 \otimes \chi_2\bigr)(\Zp)^{\mathrm{an}}_{r}$ 中稠密, 故 $\Phi$ 恒为零, 映射 $\varphi_r$ 是 $K_p(s)$-等变的.
由于 $K_p(s)$ 在 $\bigl(\pind{}_{B(\Qp)}^{\GL_2(\Qp)} \chi_1 \otimes \chi_2\bigr)(\Zp)^{\mathrm{an}}_{r}$ 上的作用是 $\mathcal{K}_p(s)$-刚性解析的, 可见 $\varphi_r$ 经一个连续 $K_p(s)$-等变映射分解:
\[
\bigl(\pind{}_{B(\Qp)}^{\GL_2(\Qp)} \chi_1 \otimes \chi_2\bigr)(\Zp)^{\mathrm{an}}_{r} \to V^{K_p(s)\text{-an}},
\]
从而特别地经一个连续 $\mathfrak{gl}_2(\Qp)$-等变映射:
\[
\bigl(\pind{}_{B(\Qp)}^{\GL_2(\Qp)} \chi_1 \otimes \chi_2\bigr)(\Zp)^{\mathrm{an}}_{r} \to V^{\mathrm{an}}.
\]
对 $r \ge r_0$ 取归纳极限并利用 \eqref{eq:36}, 得到结果.
\end{proof}

\begin{lemma}\label{lem:A.2.3}
设 $f \in \bigl(\pind{}_{B(\Qp)}^{\GL_2(\Qp)} \chi_1 \otimes \chi_2\bigr)(\Qp)^{\mathrm{an}}$ 和 $g \in \GL_2(\Qp)$ 使得 $g \cdot f$ 仍在 $\bigl(\pind{}_{B(\Qp)}^{\GL_2(\Qp)} \chi_1 \otimes \chi_2\bigr)(\Qp)^{\mathrm{an}}$ 中, 则对 $\Hom_{(\mathfrak{gl}_2(\Qp),B(\Qp))}\Bigl(\bigl(\pind{}_{B(\Qp)}^{\GL_2(\Qp)} \chi_1 \otimes \chi_2\bigr)(\Qp)^{\mathrm{an}},\ V^{\mathrm{an}}\Bigr)$ 的任何元素 $\varphi$, 有 $\varphi(g \cdot f) = g \cdot \varphi(f)$.
\end{lemma}

\begin{proof}
设 $\Omega \subset \Qp$ 为 $f$ 的支撑.
只需证明: 对一切 $z \in \Omega$ 和 $z$ 的一切紧开邻域 $\Omega_z \subset \Omega$, 有:
\begin{equation}
\label{eq:40}
\varphi\bigl(g \cdot (f|_{\Omega_z})\bigr) = g \cdot \varphi(f|_{\Omega_z}).
\end{equation}
固定 $z \in \Omega$ 和 $n := \begin{pmatrix} 1 & -z \\ 0 & 1 \end{pmatrix} \in N(\Qp)$ (于是 $z$ 是 $0 \in \Qp$ 被 $n$ 平移所得, 参见 \eqref{eq:1}).
由于 $g \cdot f$ 的支撑按假设包含在 $\Qp$ 中, $z$ 在 $g$ 下的像位于 $\Qp$ 中, 因此 $ng^{-1}$ 把 $0 \in \Qp$ 送到 $\Qp$ 中的一个点.
因此可以写成:
\[
ng^{-1} = n^{-1}b^{-1},
\]
其中 $n \in N(\Qp)$, $b \in B(\Qp)$.
由 $\varphi$ 在 $B(\Qp)$ 下的不变性, 对 $z$ 的一切紧开邻域 $\Omega_z \subset \Omega$ 得到:
\[
\varphi\bigl(g \cdot (f|_{\Omega_z})\bigr) = \varphi\bigl(bnn \cdot (f|_{\Omega_z})\bigr) = b \cdot \varphi\bigl(nn \cdot (f|_{\Omega_z})\bigr),
\]
\[
g \cdot \varphi(f|_{\Omega_z}) = bnn \cdot \varphi(f|_{\Omega_z}) = bn \cdot \varphi\bigl(n \cdot (f|_{\Omega_z})\bigr).
\]
必要时把 $f$ 换成 $n \cdot f$, $z$ 换成 $0$, 归结为证明: 对 $n \in N(\Qp)$ 和 $\Omega$ 中 $0$ 的一个足够小的邻域 $\Omega_0$, 有:
\begin{equation}
\label{eq:41}
\varphi(n \cdot f) = n \cdot \varphi(f).
\end{equation}
取 $r$ 使得 $f$ 在 $D(0,r)$ 上解析, 并取 $s \ge r$ 足够大使得 $\varphi\bigl(\bigl(\pind{}_{B(\Qp)}^{\GL_2(\Qp)} \chi_1 \otimes \chi_2\bigr)(\Qp)^{\mathrm{an}}_{r}\bigr) \subset V^{K_p(s)\text{-an}}$ (这样的 $s$ 存在, 因为 $\bigl(\pind{}_{B(\Qp)}^{\GL_2(\Qp)} \chi_1 \otimes \chi_2\bigr)(\Qp)^{\mathrm{an}}_{r}$ 是 Banach 空间而 $V^{\mathrm{an}}$ 是 Banach 空间 $V^{K_p(s)\text{-an}}$ 的并, 参见例如 \cite[Prop.~1 p.I.20]{ref:6}).
取 $t_0 \ge 0$ 使得 $\wp^{-t_0}n\wp^{t_0} \in K_p(s) \cap N(\Qp)$ (回顾 $\wp := \begin{pmatrix} p & 0 \\ 0 & 1 \end{pmatrix}$), 则对一切 $t \ge t_0$ 也有:
\begin{equation}
\label{eq:42}
\wp^{-t}n\wp^{t} \in K_p(s) \cap N(\Qp),
\end{equation}
并注意对 $t \ge 0$, $(\wp^{-t} \cdot f)|_{p^{r}\Zp} \in \bigl(\pind{}_{B(\Qp)}^{\GL_2(\Qp)} \chi_1 \otimes \chi_2\bigr)(\Qp)^{\mathrm{an}}_{r}$.
由于 $K_p(s)$ 在 $\bigl(\pind{}_{B(\Qp)}^{\GL_2(\Qp)} \chi_1 \otimes \chi_2\bigr)(\Qp)^{\mathrm{an}}_{r}$ 和 $V^{K_p(s)\text{-an}}$ 上的作用来自 $K_p(s)$ 的刚性解析作用, 且 $\varphi$ 是 $\mathfrak{gl}_2(\Qp)$-等变的, 得到:
\[
\varphi(k \cdot f') = k \cdot \varphi(f')
\]
对一切 $k \in K_p(s)$ 和 $f' \in \bigl(\pind{}_{B(\Qp)}^{\GL_2(\Qp)} \chi_1 \otimes \chi_2\bigr)(\Qp)^{\mathrm{an}}_{r}$.
于是对 $t \ge t_0$ 计算得到:
\begin{equation}
\label{eq:43}
\varphi\bigl(n \cdot f|_{p^{r+t}\Zp}\bigr)
= \wp^{t}\wp^{-t} \cdot \varphi\bigl(n\wp^{t} \cdot (\wp^{-t} \cdot f)|_{p^{r}\Zp}\bigr)
= \wp^{t} \cdot \varphi\bigl(\wp^{-t}n\wp^{t} \cdot (\wp^{-t} \cdot f)|_{p^{r}\Zp}\bigr)
\end{equation}
\[
= \wp^{t}(\wp^{-t}n\wp^{t}) \cdot \varphi\bigl((\wp^{-t} \cdot f)|_{p^{r}\Zp}\bigr)
= n \cdot \varphi\bigl(\wp^{t}(\wp^{-t} \cdot f)|_{p^{r}\Zp}\bigr)
= n \cdot \varphi\bigl(f|_{p^{r+t}\Zp}\bigr),
\]
其中第二个和第四个等号来自 $\varphi$ 的 $B(\Qp)$-等变性, 第三个来自 \eqref{eq:43} 和 \eqref{eq:42}.
对 $\Omega_0 = p^{r'}\Zp$ 且 $r' \ge r + t$, 确实得到 \eqref{eq:41}.
\end{proof}

\begin{lemma}\label{lem:A.2.4}
限制:
\begin{equation}
\label{eq:44}
\Hom_{\GL_2(\Qp)}\Bigl(\anind_{B(\Qp)}^{\GL_2(\Qp)} \chi_1 \otimes \chi_2,\ V^{\mathrm{an}}\Bigr)
\to \Hom_{(\mathfrak{gl}_2(\Qp),B(\Qp))}\Bigl(\bigl(\pind{}_{B(\Qp)}^{\GL_2(\Qp)} \chi_1 \otimes \chi_2\bigr)(\Qp)^{\mathrm{an}},\ V^{\mathrm{an}}\Bigr)
\end{equation}
是同构.
\end{lemma}

\begin{proof}
由于 $\bigl(\pind{}_{B(\Qp)}^{\GL_2(\Qp)} \chi_1 \otimes \chi_2\bigr)(\Qp)^{\mathrm{an}}$ 在 $\GL_2(\Qp)$ 作用下生成 $\anind_{B(\Qp)}^{\GL_2(\Qp)} \chi_1 \otimes \chi_2$, 映射 \eqref{eq:44} 是单射.
需要证明其满射性.
给 $L[\GL_2(\Qp)]$ 配备最细的局部凸拓扑, 考虑张量积 $L[\GL_2(\Qp)] \otimes_L \bigl(\pind{}_{B(\Qp)}^{\GL_2(\Qp)} \chi_1 \otimes \chi_2\bigr)(\Qp)^{\mathrm{an}}$ 及其归纳张量积拓扑 (此空间同构于以 $\GL_2(\Qp)$ 的元素为指标的 $\bigl(\pind{}_{B(\Qp)}^{\GL_2(\Qp)} \chi_1 \otimes \chi_2\bigr)(\Qp)^{\mathrm{an}}$ 的拷贝的直和, 该拓扑正是直和拓扑).
包含 $\bigl(\pind{}_{B(\Qp)}^{\GL_2(\Qp)} \chi_1 \otimes \chi_2\bigr)(\Qp)^{\mathrm{an}} \hookrightarrow \anind_{B(\Qp)}^{\GL_2(\Qp)} \chi_1 \otimes \chi_2$ 诱导连续 $\GL_2(\Qp)$-等变映射:
\begin{equation}
\label{eq:45}
L[\GL_2(\Qp)] \otimes_L \bigl(\pind{}_{B(\Qp)}^{\GL_2(\Qp)} \chi_1 \otimes \chi_2\bigr)(\Qp)^{\mathrm{an}}
\to \anind_{B(\Qp)}^{\GL_2(\Qp)} \chi_1 \otimes \chi_2.
\end{equation}
若 $w := \begin{pmatrix} 0 & 1 \\ 1 & 0 \end{pmatrix} \in \GL_2(\Qp)$, 有闭浸入:
\begin{equation}
\label{eq:46}
\bigl(\pind{}_{B(\Qp)}^{\GL_2(\Qp)} \chi_1 \otimes \chi_2\bigr)(\Qp)^{\mathrm{an}}
\oplus \bigl(\pind{}_{B(\Qp)}^{\GL_2(\Qp)} \chi_1 \otimes \chi_2\bigr)(\Qp)^{\mathrm{an}}
\hookrightarrow L[\GL_2(\Qp)] \otimes_L \bigl(\pind{}_{B(\Qp)}^{\GL_2(\Qp)} \chi_1 \otimes \chi_2\bigr)(\Qp)^{\mathrm{an}},
\end{equation}
定义为 $(f_1,f_2) \mapsto 1 \otimes f_1 + w \otimes f_2$.
由于 $P(\Qp) = \Qp \cup w\Qp$, \eqref{eq:45} 与 \eqref{eq:46} 的复合是连续满射.
由于它的靶与源都是紧型局部凸空间, 开映射定理蕴含它是拓扑满射.
因此 \eqref{eq:45} 也是如此.
由于 $V^{\mathrm{an}}$ 配备 $\GL_2(\Qp)$ 的连续作用, 任何连续线性映射 $\varphi : \bigl(\pind{}_{B(\Qp)}^{\GL_2(\Qp)} \chi_1 \otimes \chi_2\bigr)(\Qp)^{\mathrm{an}} \to V^{\mathrm{an}}$ 诱导连续 $\GL_2(\Qp)$-等变映射:
\[
\Phi : L[\GL_2(\Qp)] \otimes_L \bigl(\pind{}_{B(\Qp)}^{\GL_2(\Qp)} \chi_1 \otimes \chi_2\bigr)(\Qp)^{\mathrm{an}} \to V^{\mathrm{an}}.
\]
我们将证明: 若 $\varphi$ 还是 $(\mathfrak{gl}_2(\Qp),B(\Qp))$-等变的, 则 $\Phi$ 的核包含 \eqref{eq:45} 的核.
由于 \eqref{eq:45} 是拓扑满射, 这将证明 $\Phi$ 经一个连续 $\GL_2(\Qp)$-等变映射 $\overline{\Phi} : \anind_{B(\Qp)}^{\GL_2(\Qp)} \chi_1 \otimes \chi_2 \to V^{\mathrm{an}}$ 分解, 且 $\overline{\Phi}$ 限制到 $\bigl(\pind{}_{B(\Qp)}^{\GL_2(\Qp)} \chi_1 \otimes \chi_2\bigr)(\Qp)^{\mathrm{an}}$ 上与 $\varphi$ 重合 (于是 $\overline{\Phi}$ 是 $\varphi$ 的原像, 从而得到所要的满射性).
因此设 $\varphi \in \Hom_{(\mathfrak{gl}_2(\Qp),B(\Qp))}\Bigl(\bigl(\pind{}_{B(\Qp)}^{\GL_2(\Qp)} \chi_1 \otimes \chi_2\bigr)(\Qp)^{\mathrm{an}},\ V^{\mathrm{an}}\Bigr)$, 设 $g_1 \otimes f_1 + \cdots + g_r \otimes f_r$ 为 \eqref{eq:45} 的核中的元素, 其中 $g_i \in \GL_2(\Qp)$, $f_i \in \bigl(\pind{}_{B(\Qp)}^{\GL_2(\Qp)} \chi_1 \otimes \chi_2\bigr)(\Qp)^{\mathrm{an}}$:
在 $\anind_{B(\Qp)}^{\GL_2(\Qp)} \chi_1 \otimes \chi_2$ 中有:
\begin{equation}
\label{eq:47}
\sum_{i=1}^{r} g_i \cdot f_i = 0,
\end{equation}
需要证明:
\begin{equation}
\label{eq:48}
\sum_{i=1}^{r} g_i \cdot \varphi(f_i) = 0.
\end{equation}
只需证明: 若 $z \in \mathbb{P}^{1}(\Qp)$, $\Omega_z$ 是 $z$ 在 $\mathbb{P}^{1}(\Qp)$ 中一个足够小的开邻域, 则:
\begin{equation}
\label{eq:49}
\sum_{i=1}^{r} g_i \cdot \varphi\bigl(f_i|_{\Omega_z g_i^{-1}}\bigr) = 0
\end{equation}
(因为可以写 $\mathbb{P}^{1}(\Qp) = \bigcup_j \Omega_{z_j}$, $f_i = \sum_j f_i|_{\Omega_{z_j}g_i^{-1}}$, 于是确有 $\sum_i g_i \cdot \varphi(f_i) = \sum_j \sum_i g_i \cdot \varphi(f_i|_{\Omega_{z_j}g_i^{-1}}) = 0$).
设 $g := \begin{pmatrix} a & b \\ c & d \end{pmatrix} \in \GL_2(\Qp)$ 使得 $\frac{dz-b}{-cz+a} = 0$; 必要时把 $g_1 \otimes f_1 + \cdots + g_r \otimes f_r$ 换成 $g_1g^{-1} \otimes f_1 + \cdots + g_rg^{-1} \otimes f_r$, 只需处理 $z = 0$ 的情形.
若 $0$ 不在 $g_i \cdot f_i$ 的支撑中, 可以对 $\Omega_0$ 作同样假设 (必要时缩小), 则 \eqref{eq:49} 中第 $i$ 项已经为零.
因此可以假设 $0$ 在每个 $g_i \cdot f_i$ 的支撑中, $\Omega_0$ 支撑在 $\Qp$ 中, 从而每个 $g_i \cdot f_i$ 也支撑在 $\Qp$ 中 (因为我们只考虑 $f_i|_{\Omega_0 g_i^{-1}}$).
引理 \ref{lem:A.2.3} 于是结合 \eqref{eq:47} 给出:
\[
\sum_{i=1}^{r} g_i \cdot \varphi\bigl(f_i|_{\Omega_0 g_i^{-1}}\bigr)
= \sum_{i=1}^{r} \varphi\bigl(g_i \cdot (f_i|_{\Omega_0 g_i^{-1}})\bigr)
= \sum_{i=1}^{r} \varphi\bigl((g_i \cdot f_i)|_{\Omega_0}\bigr)
= \varphi\Bigl(\Bigl(\sum_{i=1}^{r} g_i \cdot f_i\Bigr)\Big|_{\Omega_0}\Bigr) = 0.
\]
\end{proof}

命题 \ref{prop:2.1.4} 由引理 \ref{lem:A.2.2} 和 \ref{lem:A.2.4} 以及同构:
\[
\Hom_{\GL_2(\Qp)}\bigl(\anind_{B(\Qp)}^{\GL_2(\Qp)} \chi_1 \otimes \chi_2,\ V^{\mathrm{an}}\bigr)
\xrightarrow{\sim} \Hom_{\GL_2(\Qp)}\bigl(\anind_{B(\Qp)}^{\GL_2(\Qp)} \chi_1 \otimes \chi_2,\ V\bigr)
\]
得出.

\begin{remark}\label{rem:A.2.5}
命题 \ref{prop:2.1.4} 的上述证明是 \cite{ref:25} 的论证在所考虑情形的具体化, 与 \cite[Prop.~2.5]{ref:21} 的证明相当类似: 引理 \ref{lem:A.2.1} 和 \ref{lem:A.2.2} 对应 \cite[Lem.~3.23]{ref:21}, 引理 \ref{lem:A.2.3} 和 \ref{lem:A.2.4} 对应 \cite[Lem.~3.1]{ref:21}.
关键区别如下: 在 \cite[Prop.~2.5]{ref:21} 的语境中, 假设 $\chi_1 \otimes \chi_2$ 是经典权 $k$ 且非临界斜率 (即 $|\chi_1(p)| > |p^{k-2}|$), 还假设 $V$ 是酉 $\GL_2(\Qp)$-Banach 空间.
这允许在引理 \ref{lem:A.2.1} 的证明中限于 $d = k-2$.
此外, 空间 $\pind{\mathrm{lp},\le k-2}_{B(\Qp)}^{\GL_2(\Qp)} \chi_1 \otimes \chi_2$ 是 $\GL_2(\Qp)$-不变的 (不只是 $(\mathfrak{gl}_2(\Qp),B(\Qp))$-不变的).
这允许陈述并证明一个避免考虑 $V$ 中局部解析向量以及 $\mathfrak{gl}_2(\Qp)$ 在 $\anind_{B(\Qp)}^{\GL_2(\Qp)} \chi_1 \otimes \chi_2$ 上作用的结果.
\end{remark}
